% =====================================================================================
% paper_sp2027/drafts/th_a_block_mixing.tex
% TH-a (SP2027_PLAN2.md Sec. 3.7, item TH (a)): block mixing at intermediate locality.
% Answers REVIEW.md, reviewer B, Weakness 5 ("Intermediate locality (1 < d < N) is not shown to
% be tight") and Question 4 ("Is the bound tight for intermediate locality d, for example
% block-wise mixing over groups of d rows?").
%
% Status: workflow draft, not \input anywhere. For the engineer's review (by Tue 3 Nov) and the
% team's hand check (TH-f, Fri 30 Oct). The file compiles on its own (tectonic); references to
% labels of the paper show as ?? here.
%
% Contents
%   Part A  body paragraph for sections/attack.tex (replaces "Tightness at intermediate locality
%           $1<d<N$ is open").
%   Part B  the other sentences to change (intro.tex twice; optional: abstract, thm:dich(b),
%           ext_attack.tex).
%   Part C  full definitions, statements and proofs for sections/appendix_attack.tex (extended
%           version).
%   Part D  compressed version for sections/appendix_short.tex (submission appendix).
%   Open points (comment at the end of the file).
%
% VERDICT. The statement holds cleanly, with one correction to the count and one caveat.
%   * Holds. Put each block of d rows through its own Reed-Solomon encoding and its own Merkle
%     tree, and run the checks of Sec. 4 block by block. A wrong claim is then accepted with
%     probability at most p^-r + prod_{i<t}(k-1-i)/(n-i), the bound of prop:amp, for EVERY d
%     (the number of blocks does not enter). Every query has support in one block, so the
%     verifier is in Spot_d(q) with q = ceil(N/d)(r+t) and qd <= (r+t)(N+d-1). With
%     beta = lambda + log2(2L) this gives 2^-lambda, within a factor
%     2(r+t)/(1-2^-lambda) = O(lambda + log L) of thm:dich(a)'s qd >= (1-eps)N, at every d.
%     These are Proposition prop:block and Corollary cor:block below.
%   * Count. The plan's q = (N/d)(1+t) holds only in two cases: the r folded rows of a block are
%     counted as ONE query (one message), or a single fold is drawn over an extension field of
%     size >= 2^beta. Under the convention of prop:amp (one query per folded row and one per
%     column, "r + t queries of support N"), the count is ceil(N/d)(r+t). The draft uses
%     prop:amp's convention.
%   * Caveat (it matters for the wording). Definition def:spot bounds the SUPPORT of an answer,
%     not its size. Opening every block whole therefore already meets thm:dich(a)'s threshold at
%     every d: q = ceil(N/d) gives zero error in the ideal-oracle model. Opening q random blocks
%     attains 1 - qd/N exactly (for d | N). Tightness of the query count at intermediate d is
%     thus immediate. What block mixing adds is SMALL ANSWERS: a folded row of k field elements,
%     or a block column of d. That lowers the price per touched row to about rk/d + t field
%     elements plus a 1/d share of the block's Merkle paths, against k weights. Mixing beats
%     opening in bytes only for d > ~4r, about 20 rows at r = 5. Remark rem:block-count states
%     this. A careful reviewer would raise it, so the body paragraph says it in one clause.
% CONDITIONS. All are stated in the proposition or in Remark rem:block-scope.
%   - C_M is honest. With a setup proof instead, k-1 becomes k-1+e for each block, and the setup
%     error sums over the block trees.
%   - No wrap-around (the range checks of Sec. 4.1).
%   - k >= 2 and n >= k, and n >= 2(k-1) for t <= beta+1. The code uses n = 4 next_pow2(k)
%     (code/src/pvi/fullcheck/commitment.py).
%   - chi is uniform and drawn after the claims.
%   - EVERY block receives the r folds and the t columns. The prover chooses which block to lie
%     about, so with t_b columns in block b the bound is p^-r + max_b (column term at t_b).
%   - One Merkle tree per block, so that an authentication path depends on its own block only.
%   - Salted leaves if the commitment must hide unopened blocks, which thm:dich(a) needs.
%   - Row layout only.
% WORDING. With this result the picture is one threshold at every locality (qd ~ N touched
%   rows), not two regimes. I recommend "lower and upper bound at every locality" over
%   "dichotomy" (REVIEW.md change 6(i); the plan's kill rule for TH (a)). Part B lists the
%   sentences to change.
%
% CHECKS DONE.
%   * Every step of the proof is elementary: Freivalds over F_p, the distance of Reed-Solomon
%     codes, sampling without replacement, and collision resistance. No external theorem is
%     used, so there was no theorem number or paper version to check on the web. Nothing below
%     cites a source outside the paper's .bib.
%   * An exhaustive computation (Python, run inline on 10 Oct 2026, not saved) found the OPTIMAL
%     cheating prover's acceptance for r = 1 on 7 toy instances:
%       - p = 13, n = 4, k = 2, t = 1, N = 2-3, with blocks {1},{2} / {1,2} / {1},{2},{3} /
%         {1,2},{3};
%       - p = 17, n = 8, k = 3, t = 1-2, N = 2, with d = 1 and d = 2.
%     The wrong claims were in the first row, the last row or every row, with shared positions.
%     The optimum never exceeds p^-1 + prod_{i<t}(k-1-i)/(n-i) (for example 0.3077 <= 0.3269
%     and 0.0924 <= 0.0945), and it is the same for every partition into blocks, as the proof
%     predicts.
%   * Numbers used:
%       - r = 5 / 7 at lambda = 128 (sections/defence.tex, "Parameters"; tab:params);
%       - 4-byte field elements in the default proof format and 31-bit ones in the compact wire
%         format (code/src/pvi/fullcheck/analytic.py, proof_bytes; claimcodec.py, field_size).
%     No measured number is used. "About 20 rows" is derived from
%     d > (bytes per field element / bytes per weight) * r.
% NOT VERIFIED / OPEN: see the "Open points" comment at the end of this file.
%
% Labels defined in Part C (extended version): prop:block, cor:block, rem:block-count,
%   rem:block-scope, eq:block-freivalds, eq:block-columns. Part D (submission appendix) defines
%   prop:block and cor:block only. The two parts go into different builds (the \ifdefined
%   switch in main.tex), so the labels do not clash.
% Labels used: def:spot, thm:dich, thm:guess, prop:amp, prop:worst, rem:merkle, rem:scope,
%   thm:sound, thm:untrusted, fact:ext, sec:defence, sec:commit, sec:intmodel, sec:protocol,
%   sec:fs, sec:related, app:untrusted. Key cited: fullversion.
% =====================================================================================

\documentclass[conference,compsoc]{IEEEtran}
\usepackage[T1]{fontenc}
\usepackage{amsmath}
\usepackage{amssymb}
\usepackage{amsthm}
\usepackage{xcolor}
\usepackage{cite}
\usepackage[hidelinks]{hyperref}
\setlength{\emergencystretch}{1.5em}

% ---- the paper's macros (main.tex), repeated so that this draft compiles on its own ----
\newcommand{\todo}[1]{{\color{red}\textbf{[TODO: #1]}}}
\newcommand{\para}[1]{\par\noindent\textbf{#1}\quad}
\newcommand{\T}{\top}
\newcommand{\F}{\mathbb{F}_p}
\newcommand{\enc}{\mathrm{Enc}}
\theoremstyle{plain}
\newtheorem{theorem}{Theorem}
\newtheorem{lemma}{Lemma}
\newtheorem{corollary}{Corollary}
\newtheorem{proposition}{Proposition}
\newtheorem{fact}{Fact}
\theoremstyle{definition}
\newtheorem{definition}{Definition}
\theoremstyle{remark}
\newtheorem{remark}{Remark}

\begin{document}

\title{TH-a Draft: Block Mixing at Every Locality}
\author{\IEEEauthorblockN{Workflow draft for review}}
\maketitle

% =====================================================================================
\section*{Part A. Body paragraph (sections/attack.tex)}
% In Sec. 3.4, paragraph "What the bounds mean": delete "Tightness at intermediate locality
% $1<d<N$ is open, and" and capitalise the rest of that sentence. The paragraph then reads as
% below, and the new paragraph "Every locality" follows it. Measured in a two-column IEEEtran
% build of this file: 8 printed lines, so net about +6 lines against the deleted clause. If
% Option B3 (generalising thm:dich(b)) is taken, the second sentence can cite
% Theorem~\ref{thm:dich}(b) instead of the proposition.

\para{What the bounds mean.} In the worst case nothing beats uniform sampling, which reaches $s/N$
on every dense network; $2^{-40}$ then takes about $27.7N$ paths per query. We do not claim that the
instance of~(a) embeds in a transformer (Remark~\ref{rem:scope}). The lower bound formalises a
folklore calculation (Section~\ref{sec:related}) for adaptive, value-aware verifiers that read the
whole trace: honest recomputation moves the change into the weights, so the model, not the trace,
must be encoded.

\para{Every locality.} Part~(a) is tight at every $d$: opening $q$ random blocks of $d$ rows catches
a forgery with probability at least $qd/N$ ($d\mid N$), but downloads rows. Running the checks of
Section~\ref{sec:defence} block by block, each encoded and committed on its own, keeps the error of
$d=N$ with $qd=O(\lambda+\log L)N$ (Proposition~\ref{prop:block}). A touched row costs one field
element per opened column plus a $1/d$ share of its block's folded rows, so mixing pays only for large
blocks.

% Longer alternative (about 12 printed lines), if the page ledger allows:
% \para{Every locality.} Part~(a) is tight at every $d$: opening $q$ random blocks of $d$ rows
% catches a forgery with probability at least $qd/N$ ($d\mid N$), but only by downloading rows.
% Mixing also works at every $d$. With each block Reed--Solomon-encoded and Merkle-committed on its
% own, the checks of Section~\ref{sec:defence}, run block by block, keep the error bound of $d=N$,
% since a lie about a block must pass that block's own columns (Proposition~\ref{prop:block}). So
% $2^{-\lambda}$ costs $qd=O(\lambda+\log L)N$, and a touched row costs one field element per
% opened column plus a $1/d$ share of its block's folded rows, not a row of weights; mixing pays
% only for large blocks.

% =====================================================================================
\section*{Part B. Other sentences to change}

\para{B1 (intro.tex, ll.~73--75).} Replace ``Conversely, symbols that mix all rows, such as
Reed--Solomon columns, reach $2^{-\lambda}$ with $O(\lambda)$ queries per layer
(Theorem~\ref{thm:dich}); intermediate locality is open.'' by:

\smallskip
\noindent Conversely, symbols that mix all rows, such as Reed--Solomon columns, reach $2^{-\lambda}$
with $O(\lambda)$ queries per layer (Theorem~\ref{thm:dich}), and symbols that mix blocks of $d$ rows
reach it with $O(\lambda)$ queries per block, so mixing matches the bound up to a factor $O(\lambda)$
at every locality (Proposition~\ref{prop:block}).
\smallskip

\para{B2 (intro.tex, l.~124, contributions).} Replace ``A lower bound for spot-checking verifiers
with a matching upper bound at full locality,'' by ``A lower bound for spot-checking verifiers,
matched up to $O(\lambda)$ at every locality,''.

\para{B3 (optional; attack.tex, Theorem~\ref{thm:dich}(b)).} Generalising part~(b) makes the
block result part of the theorem:

\smallskip
\noindent (b)~\emph{Mixing.} If each query mixes all rows of one block of $d$ rows, as the folded
rows and code columns of Section~\ref{sec:defence} do for $d=N$, then under the conditions of
Proposition~\ref{prop:block}, $O(\lambda+\log L)$ queries per block give soundness error
$2^{-\lambda}$ for a network of $L$ weight layers. For $L\le2^{\lambda}$ this is $O(\lambda)$, so
$qd=O(\lambda)N$ suffices at every $d$.
\smallskip

\noindent The proof idea then ends ``Part~(b) is our defence's soundness argument for one layer, run
block by block (Propositions~\ref{prop:amp} and~\ref{prop:block}).'' In both appendices, the proof
of Theorem~\ref{thm:dich}(b) sets $\beta=\lambda+\log_2(2L)$ in Proposition~\ref{prop:block}, which
is Proposition~\ref{prop:amp} for $d=N$, and uses $qd\le(r+t)(N+d-1)<2(r+t)N$.

\para{B4 (optional; abstract, frozen on Mon 9 Nov).} ``while commitments that mix all rows reach
$2^{-\lambda}$ with $O(\lambda)$ queries'' may become ``while commitments that mix blocks of $d$ rows
reach $2^{-\lambda}$ with $O(\lambda)$ queries per block''.

\para{B5 (optional; ext\_attack.tex, ``The two regimes'').} Replace ``The two regimes of
Theorem~\ref{thm:dich} differ by a factor $O(\lambda)$ in $qd$, and in both the verifier touches every
row. Mixing lowers the cost of each touched row from a full row of weights to $O(\lambda)$ field
elements.'' by:

\smallskip
\noindent At every locality the two bounds differ by a factor $O(\lambda+\log L)$ in $qd$, and the
verifier touches every row (Corollary~\ref{cor:block}). Mixing a block of $d$ rows lowers the cost of
each touched row from a full row of weights to one field element per opened column plus a $1/d$ share
of the block's folded rows, which is $O(\lambda)$ field elements at $d=N$ when $k\le N$.
\smallskip

% =====================================================================================
\section*{Part C. Extended version (sections/appendix\_attack.tex)}
% Insert after the paragraph "Part~(a) is the one-sided randomised query complexity of
% $\mathrm{OR}_N$ ..." that follows the proof of Theorem thm:dich, before Remark rem:merkle.
% Add to the file header: "Labels defined here: ... prop:block, cor:block, rem:block-count,
% rem:block-scope, eq:block-freivalds, eq:block-columns".

\para{Block commitments.} Let a weight layer compute $Z=A[X;\mathbf 1]$ with $A\in\F^{N\times k}$ and
$k\ge2$, where the layer's integer arithmetic embeds in $\F$ without wrap-around, as in
Proposition~\ref{prop:amp}, and let $1\le d\le N$. Partition $[N]$ into $m=\lceil N/d\rceil$
\emph{blocks} $\mathcal B_1,\dots,\mathcal B_m$ of at most $d$ rows, for instance consecutive ones.
For a matrix $Y$ with $N$ rows, $Y_b$ denotes its rows in $\mathcal B_b$. The \emph{block commitment
of locality $d$} encodes every row of $A$ with the Reed--Solomon code of Section~\ref{sec:commit}, at
the same $n\ge k$ points, which gives the encoded blocks $\hat A_b$. For each block it builds a Merkle
tree $T_b$ of its own over the $n$ columns $\hat A_{b,c}\in\F^{|\mathcal B_b|}$ of $\hat A_b$. When
unopened rows must stay hidden, leaf $c$ of $T_b$ is $H(s_{b,c}\Vert\hat A_{b,c})$ with a fresh salt
$s_{b,c}$ (Remark~\ref{rem:merkle}). $C_M$ holds the $m$ roots. For $d=N$ this is the commitment of
Section~\ref{sec:commit}.

\para{The block verifier.} $V_{\mathrm{blk}}$ knows the layer's input $X$ and receives its claims
$Z$.
(1)~After $Z$ is fixed, it draws a uniform $\chi\in\F^{N\times r}$.
(2)~It receives $u_b\in\F^{r\times k}$ for every block, honestly $u_b=\chi_b^\T A_b$, so that
$\sum_bu_b$ is the folded row $\chi^\T A$ of Section~\ref{sec:protocol}, and checks
\begin{equation}
  \chi^\T Z=\Big(\sum_{b\le m}u_b\Big)[X;\mathbf 1].
  \label{eq:block-freivalds}
\end{equation}
(3)~It draws a uniform set $C\subseteq[n]$ of $t$ distinct positions and receives, for every block,
the columns $\hat A_{b,c}$, $c\in C$, with a multiproof in $T_b$. It checks the multiproofs and
\begin{equation}
  \chi_b^\T\hat A_{b,c}=\enc(u_b)_c\qquad\text{for every }b\le m\text{ and }c\in C,
  \label{eq:block-columns}
\end{equation}
where $\enc$ encodes each row of $u_b$.

\noindent Under Fiat--Shamir the challenges are derived in the same order
(Remark~\ref{rem:block-scope}). We write $u_{b,\rho}$ for row $\rho$ of $u_b$. As
Proposition~\ref{prop:amp} does, we count one query per folded row $u_{b,\rho}$ and one per opened
column.

\begin{proposition}[block mixing]
\label{prop:block}
Assume that each tree $T_b$ binds the encoding $\hat A_b$ of $A_b$, for instance because $C_M$ is
computed honestly.
(i)~$V_{\mathrm{blk}}$ accepts the honest claims and messages. Each folded row $u_{b,\rho}$, and each
opened column $\hat A_{b,c}$ with its authentication path in $T_b$, depends only on the rows in
$\mathcal B_b$ and on salts. Hence $V_{\mathrm{blk}}\in\mathrm{Spot}_d(q)$ with $q=m(r+t)$, so
$qd\le(r+t)(N+d-1)$, and $qd=(r+t)N$ if $d$ divides $N$.
(ii)~Let the prover be arbitrary, and let $X$ and $Z\neq A[X;\mathbf 1]$ be fixed before $\chi$ is
drawn. Then $V_{\mathrm{blk}}$ accepts with probability at most
\[
  p^{-r}+\prod_{i<t}\frac{k-1-i}{n-i}\;\le\;p^{-r}+\Big(\frac{k-1}{n}\Big)^{t},
\]
plus the probability that the prover finds a collision of the hash function. This is the bound of
Proposition~\ref{prop:amp}, the case $d=N$, and it does not depend on $d$.
(iii)~With $r=\lceil\beta/\log_2p\rceil$ and $t=\lceil\beta/\log_2(n/(k-1))\rceil$, the bound
of~(ii) is at most $2^{1-\beta}$. If $n\ge2(k-1)$, then $r+t\le(1+1/\log_2p)\beta+2$.
\end{proposition}

\begin{proof}
(i)~Let the claims and messages be honest. Then $\sum_bu_b=\chi^\T A$ and $Z=A[X;\mathbf 1]$,
so~\eqref{eq:block-freivalds} holds. Encoding is linear, so
$\chi_b^\T\hat A_b=\enc(\chi_b^\T A_b)=\enc(u_b)$, and~\eqref{eq:block-columns} holds; honest
multiproofs verify. The honest $u_{b,\rho}$ is a combination of the rows in $\mathcal B_b$. Every
entry of $\hat A_{b,c}$ is the encoding of one such row at one point. A multiproof in $T_b$ consists
of hashes of other leaves of $T_b$, which are functions of columns of $\hat A_b$ and of salts. So every
query has its support in $\mathcal B_b$, a set of at most $d$ rows, and for a prover message this is
the support that Definition~\ref{def:spot} assigns to it. There are $mr$ folded rows and $mt$ opened
columns, and $md=d\lceil N/d\rceil\le N+d-1$.

(ii)~Let $D=Z-A[X;\mathbf 1]$; it is nonzero and fixed before $\chi$. Unless the prover finds a
collision, every column that passes its multiproof is the committed one, that is, the column
$\hat A_{b,c}$ of the true encoding. We assume this from now on. Let $W$ be the set of blocks $b$ with
$u_b\neq\chi_b^\T A_b$. It is fixed once $\chi$ and the messages of step~(2) are, before $C$ is drawn.

\emph{Case $W=\emptyset$.} Then $\sum_bu_b=\chi^\T A$, so~\eqref{eq:block-freivalds} reads
$\chi^\T D=0$. Pick $(i,j)$ with $D_{ij}\neq0$. For each column $\chi_\rho$ of $\chi$, entry $j$ of
$\chi_\rho^\T D$ is $\chi_{i\rho}D_{ij}+c_\rho$, where $c_\rho$ does not depend on $\chi_{i\rho}$. It
vanishes for exactly one of the $p$ values of $\chi_{i\rho}$, and the $r$ columns are independent.
Hence $\Pr[\text{accept},\,W=\emptyset]\le\Pr[\chi^\T D=0]\le p^{-r}$.

\emph{Case $W\neq\emptyset$.} Fix $\chi$ and the messages of step~(2). Let $b$ be the least element
of $W$, and let $\rho$ be a row with $u_{b,\rho}\neq(\chi_b^\T A_b)_\rho$. By linearity, row $\rho$ of
$\chi_b^\T\hat A_b$ is $\enc((\chi_b^\T A_b)_\rho)$. It and $\enc(u_{b,\rho})$ evaluate two distinct
polynomials of degree less than $k$ at $n$ distinct points. They therefore agree on a set $G$ of at
most $k-1$ positions, which is fixed before $C$ is drawn. Equation~\eqref{eq:block-columns} for block
$b$, row $\rho$ and position $c$ holds only if $c\in G$, so acceptance requires $C\subseteq G$. This
has probability $\binom{|G|}{t}/\binom nt\le\prod_{i<t}\frac{k-1-i}{n-i}$, and each factor is at
most $(k-1)/n$ because $n\ge k-1$. Averaging over $\chi$ and the messages bounds
$\Pr[\text{accept},\,W\neq\emptyset]$ by the same product.

Adding the two cases proves~(ii). The argument uses only one block of $W$, which is why $m$ does not
appear, and it holds unchanged if each tree draws its own positions. For $d=N$ it is the proof of
Proposition~\ref{prop:amp}.

(iii)~We have $r\log_2p\ge\beta$ and $t\log_2(n/(k-1))\ge\beta$, so each term is at most
$2^{-\beta}$. If $n\ge2(k-1)$, then $\log_2(n/(k-1))\ge1$, so $t\le\beta+1$, and
$r\le\beta/\log_2p+1$.
\end{proof}

\begin{corollary}[the threshold at every locality]
\label{cor:block}
Let $1\le d\le N$ and $\beta=\lambda+\log_2(2L)$. Consider a network with $L$ weight operations whose
arithmetic embeds without wrap-around, and an honest $C_M$. Run $V_{\mathrm{blk}}$ on each operation
$\ell$, with $r$ and $t_\ell$ as in Proposition~\ref{prop:block}(iii) and $n_\ell\ge2(k_\ell-1)$; the
verifier recomputes the operations without weights.
(a)~The verifier accepts an output $y\neq M(x)$ with probability at most $2^{-\lambda}$, plus that of
a hash collision. It makes $q_\ell=(r+t_\ell)\lceil N_\ell/d\rceil$ queries of locality $d$ at
operation $\ell$.
(b)~Fix a target layer of width $N$, row length $k$ and $t$ columns. Every verifier that reads this
layer through a commitment of locality $d$ that reveals nothing about unopened rows, and has
soundness error at most $\varepsilon$ on every ReLU network with a hidden layer of width $N$,
satisfies $qd\ge(1-\varepsilon)N$ (Theorem~\ref{thm:dich}(a)). The verifier of~(a) has error $2^{-\lambda}$
with $qd\le(r+t)(N+d-1)$. For every $d$ the two differ by less than a factor
$2(r+t)/(1-2^{-\lambda})$, which is $O(\lambda+\log L)$, and by at most $(r+t)/(1-2^{-\lambda})$ if
$d$ divides $N$.
\end{corollary}

\begin{proof}
(a)~If $y\neq M(x)$, some claim is wrong, since the verifier recomputes every operation without
weights. All claims are fixed before any $\chi$ is drawn. The first wrong weight operation has a
correct input, which the verifier computed from $x$ and correct claims. By
Proposition~\ref{prop:block}(ii) and~(iii), that operation passes with probability at most
$2^{1-\beta}=2^{-\lambda}/L$. A union bound over the $L$ operations, as in Theorem~\ref{thm:sound},
gives $2^{-\lambda}$. The counts are those of Proposition~\ref{prop:block}(i).

(b)~The lower bound is Theorem~\ref{thm:dich}(a), and it covers the block commitment. In the
ideal-oracle model that commitment reveals nothing beyond its answers. A Merkle realisation with
salted leaves hides the blocks it does not open (Remark~\ref{rem:merkle}), and by
Proposition~\ref{prop:block}(i) its authentication paths stay inside their block. The instance of
Proposition~\ref{prop:worst} can be made an integer network that satisfies the range bound of
Section~\ref{sec:intmodel}: take integer weights of at most 7 bits, and the logits $(1,2g)$ in place
of $(1/2,g)$, which select the same class. So the two bounds hold on the same networks. The ratio is
$(r+t)(N+d-1)/((1-2^{-\lambda})N)<2(r+t)/(1-2^{-\lambda})$, and
$r+t\le(1+1/\log_2p)\beta+2=O(\lambda+\log L)$ by Proposition~\ref{prop:block}(iii).
\end{proof}

\begin{remark}[what the count measures]
\label{rem:block-count}
Definition~\ref{def:spot} bounds the support of an answer, not its size. The threshold of
Theorem~\ref{thm:dich}(a) therefore counts touched rows.

(i)~\emph{Opening.} Suppose a query on block $b$ is answered with the block's $d$ rows, from a binding
commitment of locality $d$ such as one salted hash per block. A verifier that opens $q\le m$ distinct
uniform blocks and recomputes their claims is complete. In the ideal-oracle model it catches every
forgery of the layer with probability at least $q/m$, which is $qd/N$ if $d$ divides $N$. So the bound
$1-qd/N$ of Theorem~\ref{thm:dich}(a) is attained whenever $d$ divides $N$, as Theorem~\ref{thm:guess}
states for $d=1$. With $q=m$ the error is $0$ and $qd<N+d$ for every $d$. This is the $\lceil N/d\rceil$-query algorithm for
$\mathrm{OR}_N$ with $d$-row queries, and it downloads the layer.

(ii)~\emph{Bundling.} Sent as one message, the $r$ folded rows of a block are one query, and the
block's $t$ columns, opened with one multiproof, are another. $V_{\mathrm{blk}}$ could thus be counted
as making $2m$ queries. A single fold $\chi\in\mathbb K_s^N$ over an extension field
$\mathbb K_s\supseteq\F$ of size at least $2^\beta$, such as that of Fact~\ref{fact:ext}, makes
$m(1+t)$ queries under our convention. The proof of Proposition~\ref{prop:block} goes through over
$\mathbb K_s$, with $|\mathbb K_s|^{-1}$ in place of $p^{-r}$, since $D_{ij}\neq0$ is invertible in
$\mathbb K_s$ and the Reed--Solomon code over $\mathbb K_s$ with the same points has the same
distance. The
factor $r+t$ in Proposition~\ref{prop:block}, as in Proposition~\ref{prop:amp}, comes from counting
one query per folded row and one per column.

(iii)~\emph{Mixing changes the size of the answers.} Per layer, $V_{\mathrm{blk}}$ receives $rkm$
field elements of folded rows, $tN$ of columns, and $m$ multiproofs of at most $t\lceil\log_2n\rceil$
hashes each. For $d$ dividing $N$, this is $rk/d+t$ field elements and at most
$t\lceil\log_2n\rceil/d$ hashes per touched row, against the $k$ weights of an opened row. Take
$\varphi$-byte field elements and $\omega$-byte weights. Mixing sends fewer bytes than opening only
when $\varphi(rk/d+t)$ is below about $\omega k$, which needs $d$ above about $\varphi r/\omega$. For
31-bit field elements and int8 weights this is about $4r$, or about 20 rows at $r=5$. At $d=1$,
mixing sends more than the rows. The $t$ positions are shared, so the opened entries are those of $t$
full columns of $\hat A$, as for $d=N$, cut into $m$ pieces. Locality $d$ adds the $(m-1)rk$ field
elements of the extra folded rows and $m-1$ multiproofs.

(iv)~\emph{The ideal-oracle model.} There, a checked prover message is answered honestly, so the
error of $V_{\mathrm{blk}}$ is at most $p^{-r}$. The column checks are what make a real prover's
folded rows behave as honest answers.
\end{remark}

\begin{remark}[scope of Proposition~\ref{prop:block}]
\label{rem:block-scope}
(i)~\emph{Fiat--Shamir.} The two cases of the proof are the two rounds of the round-by-round argument
of Section~\ref{sec:fs}. Call a state doomed if $Z$ is wrong; after step~(1), if $\chi^\T D\neq0$;
and after step~(2), if $W\neq\emptyset$ or~\eqref{eq:block-freivalds} fails. Step~(1) leaves a doomed
state with probability at most $p^{-r}$, and no message of step~(2) leaves one. Step~(3) leaves one
with probability at most the column term. The Fiat--Shamir bound of Section~\ref{sec:fs} therefore
applies with the error of Proposition~\ref{prop:block}(ii).

(ii)~\emph{Untrusted commitments.} With a provider-supplied $C_M$, Theorem~\ref{thm:untrusted}
applies with each tree $T_b$ as a member of $\mathcal J$. Its setup error $\varepsilon_0$ then sums
over the $m$ trees of each matrix, so $\beta_0$ grows by $\log_2m$. Online,
Proposition~\ref{prop:block} holds for the matrix $A^*$ of Theorem~\ref{thm:untrusted}, with $k-1$
replaced by $k-1+e$ for the largest radius $e$ of the blocks, since the corrupted columns of block $b$
join $G$ (Appendix~\ref{app:untrusted}). We state Corollary~\ref{cor:block} for an honest $C_M$, as
Theorem~\ref{thm:dich} is stated. A public setup proof opens columns of every block, and we do not
claim the lower bound in its presence.

(iii)~\emph{Other codes.} For any row-wise linear code of relative distance $\delta$, the column term
becomes $(1-\delta)^t$, as in Proposition~\ref{prop:amp}.

(iv)~\emph{Spread of the queries.} Every block needs its $r$ folded rows and $t$ columns, since the
prover chooses the blocks it lies about. If block $b$ receives $t_b$ columns, the proof gives
$p^{-r}+\max_b\prod_{i<t_b}\frac{k-1-i}{n-i}$, which is no bound at all if some $t_b=0$. Sampling
blocks does not help either: by Theorem~\ref{thm:dich}(a), error $\varepsilon$ needs
$qd\ge(1-\varepsilon)N$.

(v)~\emph{Not covered.} The transposed layout of Section~\ref{sec:commit}.
\end{remark}

% =====================================================================================
\section*{Part D. Submission appendix (sections/appendix\_short.tex)}
% Insert after Proposition prop:amp and its one-line proof note, before the proof of Theorem
% thm:dich. Add prop:block and cor:block to the "Labels defined" line of the file header. If B3
% is taken, change the proof of thm:dich(b) there to "Set beta = lambda + log2(2L) in
% Proposition~\ref{prop:block} (Proposition~\ref{prop:amp} for d = N) with a union bound over the
% L layers."

\begin{proposition}[block mixing]
\label{prop:block}
In Proposition~\ref{prop:amp}, split the rows of $A$ into $m=\lceil N/d\rceil$ blocks $\mathcal B_b$
of at most $d$ rows. Reed--Solomon-encode every row at the same $n$ points, and commit the columns of
each encoded block $\hat A_b$ in a Merkle tree of its own. After $Z$, the verifier draws a uniform
$\chi\in\F^{N\times r}$ and receives $u_b$ (honestly $\chi_b^\T A_b$, where $\chi_b$ holds the rows of
$\chi$ in $\mathcal B_b$). It checks $\chi^\T Z=(\sum_bu_b)[X;\mathbf 1]$, opens the same $t$
distinct uniform positions $c$ in every tree, and checks $\chi_b^\T\hat A_{b,c}=\enc(u_b)_c$. Every
folded row and column has its support in one block, so the verifier lies in
$\mathrm{Spot}_d(m(r+t))$. With an honest $C_M$ it accepts a wrong $Z$ with probability at most
$p^{-r}+\prod_{i<t}\frac{k-1-i}{n-i}$, plus that of a hash collision, for every $d$.
\end{proposition}
\begin{proof}[Proof sketch]
Let $D=Z-A[X;\mathbf 1]\neq0$. If every $u_b$ is honest, the check reads $\chi^\T D=0$, which has
probability at most $p^{-r}$. Otherwise, take the first block with a wrong $u_b$. There, a wrong row of
$u_b$ and the same row of $\chi_b^\T\hat A_b$ are distinct codewords. They agree on at most $k-1$
positions, fixed before the positions are drawn. The full proof is in the extended
version~\cite{fullversion}.
\end{proof}

\begin{corollary}[every locality]
\label{cor:block}
Let $\beta=\lambda+\log_2(2L)$, let $r$ and $t$ be as in Proposition~\ref{prop:amp}, and let
$n\ge2(k-1)$. Then Proposition~\ref{prop:block}, with a union bound over the $L$ layers, gives error
$2^{-\lambda}$ with $qd\le(r+t)(N+d-1)$ at every $d$. This is within a factor
$2(r+t)/(1-2^{-\lambda})=O(\lambda+\log L)$ of the $(1-2^{-\lambda})N$ that
Theorem~\ref{thm:dich}(a) requires for error $2^{-\lambda}$. Opening $q$ uniform blocks whole attains
$1-qd/N$ ($d\mid N$), but each answer then holds $dk$ weights; with mixing, each answer is one folded
row or one block column.
\end{corollary}

% =====================================================================================
% OPEN POINTS (for the engineer's review and TH-f; none of them blocks the statement)
%
% 1. Counting convention. Definition def:spot does not say how a multi-vector prover message or a
%    multi-column opening is counted, and it does not bound the size of an answer. The draft
%    meets this head-on (Remark rem:block-count; one clause in the body). An alternative is one
%    sentence in def:spot: "we count one query per folded row and per opened column". The
%    lower bound does not depend on that choice.
%
% 2. Wording. After this result, "dichotomy" describes the picture less well than before. There
%    is one threshold (about N touched rows) at every locality, and locality changes only the
%    price per touched row. I recommend dropping the word (REVIEW.md change 6(i)), and keeping
%    "Theorem [lower and upper bound]" and "tight at every locality". The label sec:dichotomy can
%    stay.
%
% 3. Sum or max (REVIEW.md B-Q6; TH-e, not this item). In interactive mode C, all claims are sent
%    before any challenge, and the first wrong weight operation is fixed by them. The proof of
%    Corollary cor:block(a) therefore gives max_l (per-operation error), not the sum. The draft
%    keeps the paper's union bound, which is a correct but looser statement. The same
%    observation applies to Theorem thm:sound; TH-e should check it before using it.
%
% 4. Page budget. Part A adds about 6 printed lines net to Sec. 3.4, and B1 adds about 1 line to
%    the introduction. The page ledger in SP2027_PLAN2.md has no row for TH-a.
%
% 5. Integer hard instance. The proof of Corollary cor:block(b) changes the instance of
%    prop:worst (logits (1, 2g), weights of at most 7 bits) so that the defence's integer model
%    covers it. The authors may prefer to state prop:worst that way directly.
%
% 6. Untrusted commitments. Online soundness and the setup proof carry over (Remark
%    rem:block-scope(ii)). The LOWER bound in the presence of a public setup proof is not
%    claimed: the setup proof's opened columns are public information about every block, and
%    condition (H) has not been checked for it.
%
% 7. Not covered: the transposed layout.
%
% 8. Not addressed: a lower bound on answer SIZE. Whether a locality-d verifier needs about
%    lambda field elements per touched row (or whether the O(lambda) factor in bytes is
%    necessary) is open. The draft claims nothing about it.
%
% 9. "About 20 rows at r = 5" depends on the byte sizes and ignores the bias column, the salts
%    and the hashes. It appears only in the extended version, as an estimate.
%
% 10. TH-f: the proof of Proposition prop:block(ii) is the part to hand-check, in particular the
%     claim that W and G are fixed before C is drawn, and the use of a single block of W.
% =====================================================================================

\end{document}
